\section{Opérateur d'aire et reconstruction de l'information sectorielle}
\label{sec:area}
Le regroupement de six trits construit une bijection entre
$(\Z/9\Z)^3$ et $(\Z/27\Z)^2$, tous deux de cardinal $729$.
La bijection n'est pas présentée comme un isomorphisme additif. Sur un bloc
de bord $9\times9$, la source décompose 36 liens en deux
ensembles de 18, étiquetés par les canaux $90$ et $120$.

\subsection{Regroupement tritique et incidence du bord}
\label{sec:ii-area-incidencia-explicita}

La correspondance conserve un mot ordonné, en plus de son cardinal.
Écrivons chaque coordonnée \(x_j\in\mathbb Z/9\mathbb Z\), pour
\(j=1,2,3\), au moyen de ses représentants positionnels uniques
\(x_j=r_{2j-1}+3r_{2j}\), avec \(r_k\in\{0,1,2\}\). L'application est
\begin{equation}
 \beta_{729}(x_1,x_2,x_3)
 =\bigl(r_1+3r_2+9r_3,\ r_4+3r_5+9r_6\bigr)
 \in(\mathbb Z/27\mathbb Z)^2.
 \label{eq:ii-area-beta729}
\end{equation}

\begin{proposition}[Correspondance positionnelle entre intérieur et bord]
\label{prop:ii-area-beta729}
L'application \(\beta_{729}\) est bijective et conserve l'ordre des
six trits. Son inverse extrait les deux triplets de chiffres et les regroupe
en trois paires.
\end{proposition}
\begin{proof}
Le développement positionnel des représentants \(0,\ldots,8\) et
\(0,\ldots,26\) est unique. Extraire les six chiffres, les regrouper puis
les extraire de nouveau restitue le même mot. L'opération inverse
rétablit ainsi chacune des trois coordonnées initiales.
\end{proof}

La retenue arithmétique conserve sa position dans le mot enrichi. Les opérations
additives des deux regroupements ont des frontières différentes : par exemple,
\(\beta_{729}(0,2,0)=(18,0)\) et
\(\beta_{729}(0,1,0)=(9,0)\), dont la somme dans le tore d'ordre
vingt-sept est \((0,0)\), tandis que
\(\beta_{729}(0,3,0)=(0,1)\). Cette distinction permet de réutiliser
la correspondance positionnelle en maintenant le typage des opérations.

Dans le tore \((\mathbb Z/27\mathbb Z)^2\), le bloc
\(B_\partial=\{0,\ldots,8\}^2\) possède les ensembles de liens
orientés
\begin{align}
 \partial B_{90}
 &=\{((-1,j),(0,j)),\ ((8,j),(9,j)):0\le j\le8\},\\
 \partial B_{120}
 &=\{((i,-1),(i,0)),\ ((i,8),(i,9)):0\le i\le8\}.
 \label{eq:ii-area-enlaces}
\end{align}
Toutes les coordonnées se lisent modulo vingt-sept. Chaque famille contient
dix-huit liens et les deux sont disjointes. Leurs étiquettes conservent les
canaux d’incidence ; le cardinal des liens est \(18+18\).

\begin{theorem}[Capacité minimale de la coupe cubique]
\label{thm:ii-area-corte-minimo}
Dans le graphe cubique de sommets \(\{0,\ldots,N-1\}^3\), avec
\(N\ge2\), des liens entre voisins et une capacité unitaire par lien,
la coupe minimale entre deux faces opposées a pour capacité \(N^2\).
Pour \(N=9\), un état produit maximalement mélangé sur les
trits d'une couche minimale possède l'information
\begin{equation}
 \operatorname{cap}_{\min}=81,\qquad s_{\rm corte}=81\log3.
 \label{eq:ii-area-corte-informacion}
\end{equation}
\end{theorem}
\begin{proof}
Les \(N^2\) droites parallèles à la normale aux faces forment des chemins
disjoints par les liens. Toute coupe séparatrice intercepte chaque chemin, de
sorte qu'elle contient au moins \(N^2\) liens. Les liens traversant
une couche transversale constituent une coupe de ce cardinal. Chaque trit
uniforme apporte \(\log3\) nats ; l'entropie est additive pour l'état
produit des quatre-vingt-un trits de la couche.
\end{proof}

La capacité de la coupe cubique et l'état des trente-six liens
de \eqref{eq:ii-area-enlaces} appartiennent à des observables définies sur
des coupes différentes. La généalogie conserve les deux domaines et leurs mesures.

\subsection{Transport des modes collectifs d'incidence}
\label{sec:ii-area-transporte-colectivo}

Les espaces \(\mathcal F_6\) et \(\mathcal F_8\) de la
section~\ref{sec:barbero} fournissent les vecteurs uniformes
\[
 u_{90}=\frac1{\sqrt{90}}\sum_{f\in\mathcal F_6}(\delta_f,0),
 \qquad
 u_{120}=\frac1{\sqrt{120}}\sum_{f\in\mathcal F_8}(0,\delta_f)
\]
dans \(\ell^2(\mathcal F_6)\oplus\ell^2(\mathcal F_8)\).
Dans \(\ell^2(\partial B_{90})\oplus\ell^2(\partial B_{120})\),
on construit
\[
 v_{90}=\frac1{\sqrt{18}}\sum_{e\in\partial B_{90}}(\delta_e,0),
 \qquad
 v_{120}=\frac1{\sqrt{18}}\sum_{e\in\partial B_{120}}(0,\delta_e).
\]
Les deux paires sont orthonormales. Notons
\(\mathscr H_{\rm inc}^{\rm coll}\) et
\(\mathscr H_\partial^{\rm coll}\) les plans qu'elles engendrent respectivement.

\begin{theorem}[Entrelacement collectif des étiquettes sectorielles]
\label{thm:ii-area-entrelazamiento-colectivo}
Il existe une unique isométrie linéaire surjective entre ces plans telle que
\begin{equation}
 \mathcal J_{6\to8}u_m=v_m\quad(m=90,120).
 \label{eq:ii-area-j-colectivo}
\end{equation}
Pour les opérateurs
\[
 \begin{aligned}
 D_{\rm inc}&=90|u_{90}\rangle\langle u_{90}|
             +120|u_{120}\rangle\langle u_{120}|,\\
 D_\partial&=90|v_{90}\rangle\langle v_{90}|
             +120|v_{120}\rangle\langle v_{120}|,
 \end{aligned}
\]
on a
\begin{equation}
 \mathcal J_{6\to8}D_{\rm inc}
 =D_\partial\mathcal J_{6\to8}.
 \label{eq:ii-area-j-entrelaza}
\end{equation}
Les projecteurs sectoriels et leur combinaison orientée de coefficients
\(12:-1\) sont également transportés.
\end{theorem}
\begin{proof}
L'image d'une base fixe l'unique application linéaire. Les normes et les
produits scalaires sont conservés puisque les deux bases sont orthonormales.
L'identité \eqref{eq:ii-area-j-entrelaza} vaut sur chaque vecteur
\(u_m\), ses deux membres étant \(m v_m\). Les projecteurs
s'expriment comme \((120I-D)/30\) et \((D-90I)/30\) ; leur calcul fonctionnel
linéaire et toute combinaison des deux conservent l'entrelacement.
\end{proof}

Ce transport agit sur les deux modes uniformes et leurs étiquettes
spectrales. Les incidences individuelles et les fluctuations au sein de
chaque secteur demeurent dans leurs espaces complets. Le lien conserve
exactement la donnée collective qu'utilisera la pondération aréale.

\subsection{Caractère angulaire et opérateur d'aire}
\label{sec:ii-area-caracter}

Soit $N_e=\operatorname{diag}(0,1,2)$ sur chaque lien et
$N_{90},N_{120}$ les sommes par secteur. L'échelle étant fixée par
\[
\theta_{\rm clk}=\frac{C^*}{6}\frac{\pi}{180},\qquad
a_0=\frac{1+2\cos(10^\circ)}3\theta_{\rm clk},
\]
le facteur angulaire est le caractère réel normalisé du triplet
\begin{equation}
 \chi_{10}=\frac13\Tr\operatorname{diag}
       (1,e^{i\pi/18},e^{-i\pi/18})
 =\frac{1+2\cos(10^\circ)}3,
 \qquad a_0=\chi_{10}\theta_{\rm clk}.
 \label{eq:ii-area-caracter-triplete}
\end{equation}
La somme des phases conjuguées prouve l'égalité. Dans la branche positive
considérée, \(\theta_{\rm clk}>0\), \(\chi_{10}>0\) et
\(\gamma_{\HMT}>0\) ; ainsi, \(a_0\) et l'échelle aréale sont
positifs. Le triplet angulaire et sa normalisation précèdent la
distribution entropique sur les liens. Avec ces quantités,
l'opérateur aréal normalisé est
\[
\widehat{\mathcal A}/\ell_P^2
=\gamma_{\HMT}a_0(N_{90}+N_{120}).
\]
Son spectre prend les valeurs $n\gamma_{\HMT}a_0$, $0\le n\le72$,
avec les multiplicités données par $(1+z+z^2)^{36}$.
La preuve consiste à sommer les valeurs propres $0,1,2$ des
36 facteurs tensoriels \cite{area}.

\begin{theorem}[Récupération des secteurs de bord]
Soient
\[
N=N_{90}+N_{120},\qquad D=90N_{90}+120N_{120}.
\]
Alors
\begin{equation}
N_{90}=4N-\frac D{30},\qquad
N_{120}=\frac D{30}-3N.
\label{eq:recover}
\end{equation}
Ainsi, le total et le lecteur pondéré retrouvent conjointement
les deux occupations.
\end{theorem}
\begin{proof}
La matriz $\left(\begin{smallmatrix}1&1\\90&120\end{smallmatrix}\right)$
a pour déterminant $30$. Son inversion donne \eqref{eq:recover}.
Dans le domaine des occupations entières, l'image satisfait
$D\in30\Z$ et $90N\le D\le120N$.
\end{proof}

Le hamiltonien modulaire de la source est
$K_A=-\log\rho_A=cI+\Delta_{\rm mod}D$.
Pour $\Delta_{\rm mod}\ne0$ et une normalisation connue,
$D=(K_A-cI)/\Delta_{\rm mod}$ retrouve la seconde observable.
Cet opérateur décrit la structure spectrale de l’état réduit ;
il ne s’identifie pas par la notation au générateur d’évolution temporelle.
L’aire conserve l’occupation totale et la paire ajoute la provenance.
Par exemple, $(N_{90},N_{120})=(1,0)$ et $(0,1)$ partagent $N=1$,
mais ont $D=90$ et $D=120$.

\subsection{Normalisation de l'état et signification du Hamiltonien modulaire}
La deuxième observable est explicitée par l'état qui la définit.
La décomposition tensorielle précédente est maintenue : dix-huit qutrits
par secteur, chacun muni de l'opérateur nombre $\operatorname{diag}(0,1,2)$ ;
$N_m$ est la somme de ces opérateurs et les secteurs agissent sur des facteurs
distincts. Soit
\[
 z_m=1+e^{-m\Delta_{\rm mod}}+e^{-2m\Delta_{\rm mod}},
 \qquad m\in\{90,120\}.
\]
La famille d’états de la source a la réalisation produit
\[
 \rho_A=Z^{-1}e^{-\Delta_{\rm mod}(90N_{90}+120N_{120})},
 \qquad Z=z_{90}^{18}z_{120}^{18}.
\]
La commutativité des opérateurs nombre et la factorisation tensorielle
prouvent l'expression de $Z$. L'état est de rang plein pour
$\Delta_{\rm mod}$ réel et fini, et son logarithme spectral donne
\[
 K_A=-\log\rho_A=(\log Z)I+\Delta_{\rm mod}D.
\]
Pour $\Delta_{\rm mod}\ne0$, $K_A$ pondère différemment les excitations
des deux secteurs. La lecture conjointe avec $N$ retrouve les deux occupations
par l'équation~\eqref{eq:recover}.

La nécessité de la lecture conjointe apparaît dans deux exemples complémentaires.
Les états d'occupation $(1,0)$ et $(0,1)$ ont le même total $N=1$
et un poids modulaire différent. Inversement, $(4,0)$ et $(0,3)$ ont
le même $D=360$ et des totaux différents. La paire $(N,D)$ sépare les deux
cas. Cette récupération concerne les occupations sectorielles, et non
chaque vecteur microscopique au sein d'une dégénérescence de ces occupations.

Lorsque $\Delta_{\rm mod}=0$, l'état est proportionnel à l'identité
et $K_A$ est scalaire : cette lecture cesse de distinguer les secteurs.
Pour $\Delta_{\rm mod}\ne0$, la normalisation connue permet de retrouver
$D$ et de le composer avec l'aire. Le rôle du Hamiltonien
modulaire est ainsi défini : il caractérise l'état réduit et sa distribution d'information.
Son inclusion dans l'article est motivée par cette fonction mathématique,
tandis que le Hamiltonien d'évolution appartient au développement dynamique
de l'action.


\subsection{État réduit, décomposition de Schmidt et entropie d'intrication}
Si $\rho_A=\sum_\nu p_\nu|\nu\rangle\langle\nu|$,
la purification
\[
|\Psi\rangle=\sum_\nu\sqrt{p_\nu}\,
|\nu\rangle_A|\nu\rangle_B
\]
satisfait $\Tr_B|\Psi\rangle\langle\Psi|=\rho_A$.
Ainsi, $-\Tr(\rho_A\log\rho_A)$ est l'entropie d'intrication
de cette bipartition pure précise. L'état élargi déclare
quel système porte cette interprétation.

Pour une involution à deux feuillets possédant un vecteur propre de chaque signe,
\[
\rho_y=\frac{e^{-yR}}{2\cosh y},\qquad R^2=I,
\]
la diagonalisation fournit
\[
S(\rho_y)=\log(2\cosh y)-y\tanh y,\qquad
D(\rho_y\Vert\rho_{-y})=2y\tanh y.
\]
L'entropie est paire en $y$ ; la distance compare explicitement
les deux orientations. Cette entropie à deux niveaux ne s'identifie pas
à la capacité $81\log3$ de la coupe triadique.
